\section{Evaluation and Results}
\label{sec:evaluation}

We evaluate our approach in two stages: 
(1) specification knowledge extraction, and 
(2) fuzzing effectiveness. 
We compare against a random baseline that uniformly samples 
test parameters without adaptation or learning.\n
\subsection{Specification Knowledge Extraction}\n
The automated extraction identified 10 
protocol rules from ARINC~664 Part~7 and generated 
96 violation test cases 
(6.9 per rule on average). 
Table~\ref{tab:phase1} shows the distribution. 
This demonstrates zero-shot extraction without manual annotation.\n
\begin{table}[ht]
\centering
\caption{Specification Knowledge Extraction Results}
\label{tab:phase1}
\small
\begin{tabular}{lc}
\toprule
\textbf{Component} & \textbf{Count} \\
\midrule
Protocol rules & 10 \\
Test cases & 96 \\
Avg. cases/rule & 6.9 \\
\midrule
\multicolumn{2}{l}{\textit{Rules by category:}} \\
\quad virtual link & 1 \\
\quad sequence numbers & 1 \\
\quad filtering policing & 1 \\
\quad switching & 1 \\
\quad priority qos & 1 \\
\quad buffer management & 1 \\
\quad frame aging & 1 \\
\quad monitoring & 1 \\
\quad redundancy & 1 \\
\quad performance & 1 \\
\bottomrule
\end{tabular}
\end{table}\n
\subsection{Random Baseline Characteristics}\n
The random baseline executed 15 iterations, 
testing 5/8 VLs and 
8 distinct protocol rules. 
It achieved an average reward of 0.854, 
providing a performance baseline for comparison. 
The uniform sampling explored the test space without adaptation.\n